% SOURCE_AUTHORITY: sga5_fr_workpass.tex, lines 10517--10538 (LF-normalized unit).
% SOURCE_SHA256: 791F4EFFC5E02832D5D77ED03518C8156D6F07E4C8238B03545DB93D883FBB28
\begin{statement}{Corolário 7.4}
Sejam $X$ um $k$-esquema projetivo, liso, conexo, de dimensão $n$, e $Y$ uma seção hiperplana lisa de $X$, definida por uma seção de um $\mathcal O_X$-módulo inversível amplo $L$. Denotando por $p$ o morfismo canônico de $X$ em $k$ e por $p_*$ o morfismo de Gysin correspondente, tem-se a igualdade
\[
\operatorname{EP}(X)-\operatorname{EP}(Y)=p_*\bigl(c(\check\Omega_X^1)c(L)^{-1}\bigr)
\tag{7.4.1}
\]
em $H^0(k,\mathbb Q_\ell)=\mathbb Q_\ell$. Além disso,
\[
b^{n-1}(Y)=2b^{\,n-1}(X)-b^n(X)+(-1)^n p_*\bigl(c(\check\Omega_X^1)c(L)^{-1}\bigr).
\tag{7.4.2}
\]
\end{statement}

\emph{Demonstração.} A igualdade (7.4.1) procede de 7.3 e de 4.2, tendo em conta a igualdade evidente
\[
1-c(L)^{-1}c_1(L)=c(L)^{-1}.
\]
Utilizando (7.2), vê-se que
\[
b^{n-1}(Y)=b^{n-1}(X)-b^n(X)+b^{n+1}(X)+(-1)^n(\operatorname{EP}(X)-\operatorname{EP}(Y)).
\]
Mas $b^{n-1}(X)=b^{n+1}(X)$ por dualidade de Poincaré, de onde deduz-se (7.4.2).
